\section{\kcode{groupprivate} Directive}
\label{sec:groupprivate}
\index{directives!groupprivate@\kcode{groupprivate}}
\index{groupprivate directive@\kcode{groupprivate} directive}

The \kcode{groupprivate} directive introduced in Specification 6.0
allows specified list items to be replicated such that each contention 
group will have its own uninitialized copy. The list item is shared among
threads of the contention group and does not exist 
outside the scope of the contention group. 

In the following example, the variable \ucode{x} is defined as a static 
variable and specified with the \kcode{groupprivate} data attribute in 
the function \ucode{foo}. Four teams created by the \kcode{teams} construct
execute the \kcode{parallel} region that calls the 
\ucode{foo} function. For each team the groupprivate variable \ucode{x} 
is created and is accessible for the group of tasks of the \kcode{parallel} region.

\cppexample[6.0]{groupprivate}{1}

\ffreeexample[6.0]{groupprivate}{1}
